\documentclass{beamer}

\usepackage{verbatim}
\usepackage{xcolor} % See documentation PDF at http://www.ctan.org/pkg/xcolor
\definecolor{darkgreen}{rgb}{0,0.3,0}
\usepackage{tikzsymbols}


\setbeamertemplate{section in toc}[sections numbered]
\setbeamertemplate{subsection in toc}[subsections numbered]
\setbeamertemplate{subsubsection in toc}[subsubsections numbered]
\usetheme{Singapore}

\newcommand{\code}[1]{{\color{darkgreen}\texttt{#1}}}


\begin{document}

\title{Formal Models of Language \\[1.5em]
 \includegraphics[width=0.5\textwidth]{images/kitten_string_flickr_albaraa.jpg} \\[-1.0em]
 %\small{Possibilities} \\[1.0em]
 %LT1 \\[1.0em]
 }
\author{\href{http://jon.dehdari.org}{Jon Dehdari}}
\frame{\titlepage}

\section{Overview}
\begin{frame}{Introduction}
Hi
\end{frame}

% Formal language
\begin{frame}{Formal Languages}
\begin{block}{}
\begin{itemize}
	\item Once Upon a Time...
	\pause
	\item Mathematicians started to think about language...
	\pause
	\item They used ideas from logic to represent linguistic objects...
	\pause
	\item They had a craaaazy idea... \\
	\pause
	\begin{center}
	\includegraphics[width=0.6\textwidth]{images/doc_brown_strings.jpg}
	\end{center}
\end{itemize}
\end{block}
\end{frame}


\begin{frame}{Strings}
\begin{block}{What's a String?}
\begin{itemize}
	\item A \textbf{string} in this context is just a sequence of letters
	\pause
	\item A \textbf{formal language} ($L$) is a subset of all the possible strings
	\pause
	\item An \textbf{alphabet} ($\Sigma$) here is a set of all the letters in the language
	\pause
	\item Alphabets and letters don't need to correspond to writing systems used for natural languages
	\pause
	\item For example, this set:\\
	\begin{center}
	{\LARGE \{ \Smiley[][cyan], \Springtree, \Snowman \}}\\
	\end{center}
	is a perfectly valid alphabet for a formal language. \pause But we usually use boring symbols like \{a, b, c\}
\end{itemize}
\end{block}
\end{frame}


% Formal grammar: grammar, automaton
\begin{frame}{Formal Grammar}
\begin{block}{}
\begin{itemize}
	\item A formal grammar is a way of telling what a valid string is in a formal language
	\item Formal grammars can also \textbf{generate} valid strings
	\pause
% weak & strong generative capacity
	\item If two different grammars can generate \& accept the same formal languages, then they have the same \textbf{weak generative capacity}
	\pause
	\item If two different grammars can generate \& accept the same structures as well, then they have the same \textbf{strong generative capacity}
\end{itemize}
\end{block}
\end{frame}


% Chomsky Hierarchy
\begin{frame}{Formal Language Hierarchy}
\begin{center}
\begin{tabular}{|ll|}
\hline
& Formal Language \\
\hline
\hline
& Non-Turing-acceptable\\
\hline
0: & Recursively enumerable \\
\hline
& Recursive/\,Decidable \\
\hline
1: & Context-sensitive \\
\hline
& Indexed \\
\hline
& Mildly context-sensitive \\
\hline
2: & Context-free \\
\hline
& Deterministic context-free \\
\hline
3: & Regular \\
\hline
& Finite \\
\hline
\end{tabular}
\end{center}
\pause
\footnotesize{This is extended from the older \textit{Chomsky hierarchy}.}
\end{frame}


% applications of knowing this stuff
\begin{frame}{Why is this Stuff Relevant??}
\begin{block}{}
\begin{itemize}
	\item Knowing what types of formal languages a grammar/automaton can generate \& accept will give you an idea of what phenomena in natural languages that they can handle
	\pause
	\item For example: long-distance dependencies, complex reordering in machine translation, reduplication, etc.
	\pause
	\item You can also get an idea of how fast or slow it will take for a computer (or human) to process sequential stuff (like natural language!)
\end{itemize}
\end{block}
\end{frame}


\begin{frame}{Relevant Formal Languages for Dealing with Natural Languages}
\begin{block}{}
\begin{itemize}
	\item 
	\item 
	\item 
\end{itemize}
\end{block}
\end{frame}


\begin{frame}{Finite Languages}
\begin{block}{}
\begin{itemize}
	\item All of the 
	\item Time: constant (through hash-table lookup)
	\item Memory: constant (duh)
\end{itemize}
\end{block}
\end{frame}


% Regular language: fixed history: linear time, constant memory.  Regular grammars can actually do alot
\begin{frame}{}
\begin{block}{}
\begin{itemize}
	\item 
	\item 
	\item 
\end{itemize}
\end{block}
\end{frame}


% DCF lang: full history, as long as unambiguous: linear time, log^2(n) memory
\begin{frame}{}
\begin{block}{}
\begin{itemize}
	\item 
	\item 
	\item 
\end{itemize}
\end{block}
\end{frame}


% CF lang: full history, can be ambiguous, allow center-embedding: about n^3 time (basic), ?? mem
\begin{frame}{}
\begin{block}{}
\begin{itemize}
	\item 
	\item 
	\item 
\end{itemize}
\end{block}
\end{frame}


% MCS lang: same, allows full reduplication (cross-serial deps) and a^n b^n c^n d^n: n^6 time. TAG, CCG, LIG, Head Grammars, 
\begin{frame}{}
\begin{block}{}
\begin{itemize}
	\item 
	\item 
	\item 
\end{itemize}
\end{block}
\end{frame}


% Recursively enumerable langs: allows any string that a computer (eg. Turing machine, untyped \lambda calculus, MLP) can generate
\begin{frame}{}
\begin{block}{}
\begin{itemize}
	\item 
	\item 
	\item 
\end{itemize}
\end{block}
\end{frame}



\begin{frame}{}
\begin{block}{}
\begin{itemize}
	\item 
	\item 
	\item 
\end{itemize}
\end{block}
\end{frame}


\end{document}
